\documentclass[10pt,a4paper]{article}
\usepackage[T1]{fontenc}
\usepackage{lmodern}
\usepackage{amsmath,amssymb,amsthm,mathtools}
\usepackage[margin=0.77in]{geometry}
\usepackage{microtype}
\usepackage[colorlinks=true,linkcolor=blue,citecolor=blue,urlcolor=blue]{hyperref}
\setlength{\parindent}{0pt}
\setlength{\parskip}{4pt}
\setlength{\abovedisplayskip}{7pt}
\setlength{\belowdisplayskip}{7pt}
\newtheorem{proposition}{Proposition}[section]
\newcommand{\R}{\mathbb R}
\newcommand{\norm}[1]{\left\|#1\right\|}
\newcommand{\ip}[2]{\left\langle#1,#2\right\rangle}
\newcommand{\pos}[1]{\left[#1\right]_+}
\newcommand{\simplex}{\Delta_m}
\title{Mathematical notes on executable common ascent,\protect\\
cross-view invariants, and decision-space uncertainty}
\author{RefCloth exploration notes}
\date{31 October 2026}
\begin{document}
\maketitle
\vspace{-8pt}

\textbf{Scope.} These notes specify three falsifiable source-side candidates.
Their ingredients---linear-programming duality, multiobjective ascent,
factorial/M\"obius contrasts, and ellipsoidal robust optimization---are standard
mathematics, not claimed new general theorems. The application hypotheses must
be distinguished from these identities. No experimental outcomes are reported
here; no result below establishes task failure, closed-loop transfer, or an
unconditional deployment certificate.

\section{Common ascent in the final executable box}
Fix a texture-parameter vector $\theta\in[-b,b]^p$, step budget $\eta>0$, and
source task--state objectives $J_i$, $i=1,\ldots,m$. Put
$g_i=\nabla J_i(\theta)$ and $q_i=g_i/s_i$, where $s_i>0$ is fixed for this
step. Raw gradients use $s_i=1$; another choice is
$s_i=\max(\norm{g_i}_1,\varepsilon)$. Zero-gradient groups must be retained.
Common regularizers belong in every $J_i$ if their gradients enter the update.

An update followed by parameter clipping can execute precisely the increments
\begin{equation}\label{eq:box}
 B=\{\delta:\ell_j\leq\delta_j\leq u_j\},\quad
 \ell_j=\max\{-\eta,-b-\theta_j\},\qquad
 u_j=\min\{\eta,b-\theta_j\}.
\end{equation}
This box contains zero but need not be symmetric. With fixed offsets $c_i$, solve
\begin{equation}\label{eq:lp}
 \gamma^*=\max_{\delta\in B}\min_i(q_i^\top\delta+c_i)
 =\max_{\delta,t}\{t:q_i^\top\delta+c_i\geq t\ \forall i,\ \delta\in B\}.
\end{equation}
The executable output is the \emph{primal increment} $\delta^*$ itself.
Normalization and execution geometry are distinct design choices.

\begin{proposition}[Support-function dual]
For $\simplex=\{w\geq0:\sum_iw_i=1\}$ and $h(w)=\sum_iw_iq_i$,
\begin{equation}\label{eq:dual}
 \gamma^*=\min_{w\in\simplex}
 \left\{\sum_{j=1}^p\max\{\ell_jh_j(w),u_jh_j(w)\}+w^\top c\right\}.
\end{equation}
Every feasible $\delta$ gives a lower bound in \eqref{eq:lp}; every simplex
weight vector gives an upper bound in \eqref{eq:dual}. Their difference is a
directly checkable primal--dual gap.
\end{proposition}
\begin{proof}
Replace the finite minimum over groups by the minimum over simplex mixtures.
The resulting objective is affine in each variable on convex compact sets,
so minimax interchange applies. Maximizing $h^\top\delta$ over $B$ separates
coordinate-wise, giving the support function displayed in \eqref{eq:dual}.
\end{proof}

For the full symmetric cube $B=[-\eta,\eta]^p$ and $c=0$,
\begin{equation}\label{eq:l1dual}
 \gamma^*=\eta\min_{w\in\simplex}\norm{\textstyle\sum_iw_iq_i}_1.
\end{equation}
Thus $\gamma^*=0$ exactly when $0\in\operatorname{conv}\{q_i\}$ in this
symmetric case. More generally, with $c=0$, $\gamma^*>0$ is equivalent to
existence of a strict common-ascent increment \emph{in the actual box}.
The convex-hull equivalence cannot be transferred without qualification to a
clipped box with zero on its boundary.

Multiobjective gradient combination predates this application
\cite{sener,pcgrad}. What must be tested is whether respecting the executed
box changes native updates, beyond using existing gradient combination alone.

\newpage
\subsection*{Why a final sign can invalidate a common-ascent direction}
Take two linear objectives with $g_1=(1,0)$, $g_2=(-2,1)$ and
$B=[-1,1]^2$. The sign of their mean is $(-1,1)$, with objective gains
$(-1,3)$. Projecting each conflicting gradient against the original other
gradient, as in the two-group PCGrad calculation \cite{pcgrad}, gives
\[
 \widetilde g_1=(1/5,2/5),\qquad \widetilde g_2=(0,1).
\]
Their sum has positive gains $(1/5,1)$, whereas its sign $(1,1)$ has gains
$(1,-1)$. A strict common-ascent direction existed and was destroyed by signing.
The maximin primal solution instead is
\[
 \delta^*=(1/3,1),\qquad \gamma^*=1/3.
\]
Indeed, gains at least $t$ require $\delta_1\geq t$ and
$\delta_2\geq t+2\delta_1\geq3t$, hence $t\leq1/3$.
Even signing the optimal dual combination is insufficient: weights
$(2/3,1/3)$ give $h=(0,1/3)$, whose sign $(0,1)$ has worst gain zero.
The missing fractional coordinate is determined by active primal constraints.
For $\ell_1$-normalized rows $(1,0)$ and $(-2/3,1/3)$, the sign of their mean
still harms the second group; the normalized LP optimum is $(1/5,1)$.

\begin{proposition}[Sparse fractionality of an optimal basic solution]
After eliminating fixed coordinates of $B$, there exists an optimal basic
solution with at most $m-1$ increment coordinates strictly inside their bounds.
If that solution has $k$ active group constraints, the sharper limit is $k-1$.
\end{proposition}
\begin{proof}
There are $p+1$ free LP variables $(\delta,t)$. If $r$ increment coordinates
are interior, at most $p-r$ independent box constraints are active. A basic
solution requires $p+1$ independent active constraints, so
$p-r+k\geq p+1$. An optimal vertex exists for this pointed LP with bounded
increment box and finite optimum. The claim concerns an optimal basic
solution, not every point on its optimal face or an averaged solver iterate.
\end{proof}

\begin{proposition}[Conditional native-export gain bound]\label{prop:native}
Let $F_i(\theta)=\mathcal F_i(T_\theta)$ be a continuously relaxed native
objective. At the current point and along the relevant box segment, assume
\[
 \norm{\nabla F_i(\theta)-g_i}_1\leq\rho_i,\qquad
 \norm{\nabla F_i(a)-\nabla F_i(a')}_1\leq L_i\norm{a-a'}_\infty.
\]
Suppose $\mathcal F_i$ is $\kappa_i$-Lipschitz in texture $\ell_\infty$ norm,
and the actual export $Q$ satisfies $\norm{Q(T)-T}_\infty\leq q$. Then
\begin{equation}\label{eq:nativebound}
 \mathcal F_i(Q(T_{\theta+\delta}))-\mathcal F_i(Q(T_\theta))
 \geq g_i^\top\delta-\rho_i\eta-\tfrac12L_i\eta^2-2\kappa_iq.
\end{equation}
\end{proposition}
\begin{proof}
Integrate the native gradient along the step. H\"older's inequality bounds its
initial discrepancy by $\rho_i\eta$, and the integrated remainder by
$L_i\eta^2/2$. Comparing quantized and relaxed textures at both endpoints costs
at most $2\kappa_iq$.
\end{proof}
Dividing $F_i,g_i$ and the remainder by the same fixed $s_i$, the LP offset
$c_i=-(\rho_i\eta+L_i\eta^2/2+2\kappa_iq)/s_i$ gives a conditional common
normalized native-gain lower bound $\gamma^*$.

These bounds require justified gradient-error, curvature, and export constants.
Finite probes do not establish uniform constants. Nearest 8-bit rounding has
$q\leq1/(2\cdot255)$; truncation may require $q\leq1/255$. Identical exported
PNGs have zero native gain despite a positive relaxed prediction. Clipping and
nonsmooth TV also require their own smoothness qualifications.

\newpage
\section{Four-corner interaction as a conditional rendering invariant}
Fix simulator state, instruction, proprioception, inference noise, and action
coordinates. Let $A_P(u,v)\in\R^d$ denote the selected action vector when the
two views use textures $u,v$ through path $P\in\{R,N\}$ (synthesis/native).
Physical deployment uses the shared diagonal $u=v=\theta$; off-diagonal
evaluations are camera-input interventions, not separately deployable textures.
Put $A_{P,ij}=A_P(i\theta,j\theta)$ and
\begin{align}
 d_{1,P}&=A_{P,10}-A_{P,00},&
 d_{2,P}&=A_{P,01}-A_{P,00},\notag\\
 d_{J,P}&=A_{P,11}-A_{P,00},&
 I_P&=A_{P,11}-A_{P,10}-A_{P,01}+A_{P,00}.\label{eq:corners}
\end{align}
All four corners within each contrast use the same rendering path.

\begin{proposition}[Unique normalized linear invariant]
A linear contrast $\sum_{i,j\in\{0,1\}}c_{ij}A_{ij}$ with $c_{11}=1$ is
invariant to every additive nuisance $b+e_1(i)+e_2(j)$ if and only if
\begin{equation}\label{eq:unique}
 (c_{00},c_{10},c_{01},c_{11})=(1,-1,-1,1).
\end{equation}
\end{proposition}
\begin{proof}
Annihilating the first-view indicator forces $c_{10}+c_{11}=0$; annihilating
the second forces $c_{01}+c_{11}=0$; annihilating constants forces
$\sum c_{ij}=0$. These equations give \eqref{eq:unique}, which also suffices.
The additive nuisance subspace has dimension three in the four-dimensional
corner space, leaving one quotient coordinate. This is a standard
factorial/M\"obius contrast, not a new general invariance theorem \cite{rota}.
\end{proof}

Define the cross-path action error $E(u,v)=A_R(u,v)-A_N(u,v)$. Exactly,
\begin{equation}\label{eq:errorcontrast}
 I_R-I_N=E(\theta,\theta)-E(\theta,0)-E(0,\theta)+E(0,0).
\end{equation}
Thus the error cancels if $E(u,v)=b+e_1(u)+e_2(v)$ at these corners, even if
$d_{J,R}-d_{J,N}$ is large. A nonlinear fusion error need not be separable.
Replacing only $A_{R,00}$ by $A_{N,00}$ changes the contrast to
$I_R-E(0,0)$ and spoils the cancellation of a constant cross-path offset.

\begin{proposition}[Behavioral relevance requires unilateral control]
For any norm,
\begin{equation}\label{eq:jointbound}
 \norm{d_{J,N}}\geq\pos{\norm{I_N}-\norm{d_{1,N}}-\norm{d_{2,N}}}.
\end{equation}
If $\norm{I_R-I_N}\leq\xi$ and $\norm{d_{v,N}}\leq\tau_v$, then
$\norm{d_{J,N}}\geq\pos{\norm{I_R}-\xi-\tau_1-\tau_2}$.
\end{proposition}
\begin{proof}
Use $d_{J,N}=I_N+d_{1,N}+d_{2,N}$ and the reverse triangle inequality twice.
\end{proof}
RMS on a fixed $d$-coordinate subset is a scaled Euclidean norm, so the same
bound applies. Constraints on synthesis unilateral effects alone do not give
native unilateral control without an additional cross-path error bound.

\textbf{Identification limit.} Even knowing the native clean corner does not
identify native total joint displacement under unrestricted additive error.
For any action vector $z$, $A_N(i,j)=A_R(i,j)-iz$ preserves the clean corner
and interaction while changing native joint displacement by $-z$.
Consequently an invariant interaction is not a surrogate for the entire native
action effect without additional calibration or structural assumptions.

\newpage
\subsection*{Mixed curvature, relative error, and protected local updates}
For an ideal $C^2$ extension on the rectangle of texture interpolations, two
applications of the fundamental theorem of calculus give
\begin{equation}\label{eq:mixed}
 I_P(\theta)=\int_0^1\!\int_0^1
 D_{12}A_P(s\theta,t\theta)[\theta,\theta],dt\,ds.
\end{equation}
If the bilinear operator norm of $D_{12}E$, for the chosen input and output
norms, is uniformly at most $L_{12}$ there,
\begin{equation}\label{eq:mixederror}
 \norm{I_R-I_N}\leq L_{12}\norm{\theta}^2.
\end{equation}
This is an \emph{absolute} second-order error bound. The interaction itself is
generally second-order, so relative fidelity need not improve as the texture
shrinks. For example, scalar maps $A_N(u,v)=uv$ and
$E(u,v)=a uv$ give $I_N(\theta)=\theta^2$ and
$|I_R-I_N|/|I_N|=|a|$ for every nonzero $\theta$.
When $I_N=0$, relative error is undefined. A fixed denominator floor must be
specified before inspection if relative errors are summarized.

\textbf{Smoothness limit.} Quantized PNGs, finite-precision inference, clipping,
and visibility changes do not automatically admit this $C^2$ extension.
The four-corner algebra, uniqueness, and triangle inequalities remain exact
for observed outputs; the smooth mixed-curvature interpretation is conditional.
Uncoupled inference noises would additionally confound the four-corner contrast.

A possible protected source objective is
\begin{equation}\label{eq:protected}
 \max_{\theta}\ \mathbb E\norm{I_R(\theta)}^2,
 \qquad C_v(\theta)=\mathbb E\norm{d_{v,R}(\theta)}^2\leq\tau_v^2,
 \qquad T_\theta\in\mathcal T_\epsilon.
\end{equation}
An average constraint does not imply each-state protection; a worst-state or
tail claim requires the corresponding prespecified constraint. Native
unilateral effects must be tested separately, as \eqref{eq:jointbound} requires.
For $g=\nabla\norm{I_R}^2$ and $h_v=\nabla C_v$, an economical local problem is
\begin{equation}\label{eq:cone}
 \max_{\delta\in\mathcal C,\,\norm{\delta}_2\leq r}g^\top\delta,
 \qquad\mathcal C=\{\delta:h_v^\top\delta\leq0,\ v=1,2\}.
\end{equation}
Its solution is $r\Pi_{\mathcal C}(g)/\norm{\Pi_{\mathcal C}(g)}$ when the
projection is nonzero, and zero otherwise. Moreau's cone decomposition gives
the corresponding support-function identity. This is standard constrained
gradient projection, related to GEM \cite{gem}; the two-constraint dual can be
solved by enumerating active sets. If $\nabla C_v$ is $\beta_v$-Lipschitz,
$C_v(\theta+\delta)\leq C_v(\theta)+\beta_v r^2/2$ for such a first-order
feasible step. Finite-step protection still needs a margin or an actual
backtracking check.

\textbf{Local spectral observation.}
For a fixed action direction $a$, let
$H_a=\operatorname{sym}(D_{12}(a^\top A)(0,0))$. With bounded third derivatives,
\[
 a^\top I(tv)=t^2v^\top H_av+O(t^3).
\]
If $P$ projects onto $\ker D_1A(0,0)\cap\ker D_2A(0,0)$, maximizing this
signed quadratic over unit vectors in that subspace is a restricted eigenvalue
problem. Maximizing the full vector norm is instead quartic and is not solved
by one such eigenvector. First-order unilateral nulling does not remove
second-order solo changes. Also, $\norm{I(\theta)}^2=O(\norm{\theta}^4)$ near
zero, so its gradient can vanish there; initialization and optimization gains
must be separated.

\newpage
\section{Decision-space ellipsoidal renderer uncertainty}
Fix a source action head $H$ and a linear map $P$ selecting execution-relevant
coordinates in one common control-unit system. Let $z_{N,0}$ be native clean
features, $d_R=z_R(\theta)-z_R(0)$ and $d_N=z_N(\theta)-z_N(0)$. Define
\begin{equation}\label{eq:actiondisp}
 u_R=P\{H(z_{N,0}+d_R)-H(z_{N,0})\},\quad
 u_N=P\{H(z_{N,0}+d_N)-H(z_{N,0})\},\quad e=u_N-u_R.
\end{equation}
The action head is evaluated as given, without a linear approximation. The
native clean anchor is fixed. Production affine action-unnormalization scales
belong in $P$; their offsets cancel in displacements. Mixing normalized and
unnormalized coordinates would change the Euclidean objective and its geometry.

Given a fitted $C\succ0$, consider the conservative action-energy objective
\begin{equation}\label{eq:ellipsoid}
 V(u)=\min_e\{\norm{u+e}_2^2:e^\top C^{-1}e\leq r(u)^2\}.
\end{equation}
Absolute uncertainty uses $r(u)=\rho$. Relative uncertainty uses
$r(u)=\rho\norm{u}_2$, with a separately estimated dimensionless shape.
The following is a standard ellipsoidal quadratic subproblem
\cite{more,elghaoui}; its exact solution is useful for specifying the gradient.

\begin{proposition}[Inner optimum and envelope derivative]\label{prop:ellipsoid}
Write $C=U\operatorname{diag}(c_j)U^\top$, $v=U^\top u$, and suppose $r(u)>0$.
If $u^\top C^{-1}u\leq r(u)^2$, then $V(u)=0$ and $e^*=-u$.
Otherwise there is a unique $\mu>0$ satisfying
\begin{equation}\label{eq:secular}
 \sum_j\frac{c_jv_j^2}{(c_j+\mu)^2}=r(u)^2,
\end{equation}
and the worst-case remaining action and objective are
\begin{equation}\label{eq:worst}
 y^*=u+e^*=U\operatorname{diag}\!\left(\frac{\mu}{c_j+\mu}\right)v,
 \qquad V(u)=\norm{y^*}_2^2.
\end{equation}
At an exterior point, the two envelope gradients are
\begin{equation}\label{eq:envelope}
 \nabla V_{\rm abs}(u)=2y^*,\qquad
 \nabla V_{\rm rel}(u)=2\{y^*-\mu\rho^2u\}.
\end{equation}
\end{proposition}
\begin{proof}
The convex inner problem has a unique optimizer. Its Lagrangian is
$\norm{u+e}^2+\mu(e^\top C^{-1}e-r(u)^2)$. Exterior points satisfy Slater's
condition and have an active constraint. Stationarity gives
$e^*=-C(C+\mu I)^{-1}u$, yielding \eqref{eq:secular}--\eqref{eq:worst}.
The left side of \eqref{eq:secular} decreases continuously from a value above
$r(u)^2$ to zero, hence has a unique positive root. Differentiating the
Lagrangian with respect to $u$ at its optimum gives
$2y^*-\mu\nabla(r(u)^2)$, which proves \eqref{eq:envelope}.
\end{proof}

At strict interior points, both value and gradient are zero. With zero radius,
the limiting problem is simply $V(u)=\norm{u}^2$. Relative uncertainty is
positively homogeneous: $V_{\rm rel}(a u)=a^2 V_{\rm rel}(u)$ for $a\geq0$.
Omitting the second term in its envelope gradient gives the right forward
score but the wrong update. A single isotropic subtraction
$(\norm{u}-\rho)_+^2$ is only a monotone radial transformation of action energy;
on one sample its nonzero sign-gradient direction is unchanged. Anisotropic
uncertainty must be distinguished from such a scalar reweighting.

\newpage
\subsection*{What calibration and a conditional lower bound do not prove}
If the \emph{actual} error $u_N-u_R$ belongs to the chosen uncertainty set,
\begin{equation}\label{eq:conditionalvalue}
 \norm{u_N}^2\geq V(u_R),\qquad 0\leq V(u_R)\leq\norm{u_R}^2.
\end{equation}
This follows because the actual error and zero error are feasible competitors
for the inner minimum. If the new texture obeys that membership condition and
its lower bound exceeds a measured old native action energy, an improvement
of that native action energy follows. Neither an increase of the lower bound
alone nor membership measured at a different texture establishes the premise.

A source-only estimation recipe can use disjoint \emph{state groups} for shape
and radius estimation. For absolute errors, one possible origin-centered shape
is $C=\mathbb E[ee^\top]+\lambda I$; centering would remove a systematic error
that the uncertainty model may need to cover. Relative errors can instead fit
the shape of $e/\norm{u_R}$, with a declared treatment of zero or negligible
denominators. A held-out empirical quantile of Mahalanobis lengths defines an
empirical radius, not a uniform error bound. Correlated probes and state shift
further limit a naive probabilistic interpretation. Optimization-selected
textures need not follow the calibration distribution, so random-probe
coverage cannot certify them.

For a fixed comparison, signal-versus-error covariance can also define
generalized directions $S a=\lambda C a$, with $S=\mathbb E[u_Nu_N^\top]$.
That is an established generalized-eigenvalue construction. Direct action
objectives are already used in VLA texture attacks \cite{unitexture}. Thus a
meaningful distinction must be a tested advantage of the uncertainty geometry
over both direct action energy and a fixed action-space quadratic metric.

\textbf{Evaluation requirements implied by the mathematics.}
An ellipsoid can contain $-u_R$, making the robust score and its gradient zero;
this is a substantive failure of optimization usability. Flat-score states
must not silently disappear from rank summaries, nor may arbitrary argmax
ties be interpreted as useful selection. Report tie rules, degeneracy, and
each-state behavior. All candidate updates must be evaluated in the same
unweighted native control coordinates. Differing loss scales must not silently
change regularization strength. Finally, native value fidelity and native
gradient fidelity are different questions: a smaller prediction error does
not by itself justify a better executable update.

\begin{thebibliography}{9}\small
\setlength{\itemsep}{3pt}
\bibitem{sener} O. Sener and V. Koltun. \emph{Multi-Task Learning as Multi-Objective
Optimization}. 2018. \href{https://arxiv.org/abs/1810.04650}{Primary paper}.
\bibitem{pcgrad} T. Yu et al. \emph{Gradient Surgery for Multi-Task Learning}.
2020. \href{https://arxiv.org/abs/2001.06782}{Primary paper}.
\bibitem{rota} G.-C. Rota. \emph{On the Foundations of Combinatorial Theory I:
Theory of M\"obius Functions}. 1964.
\href{https://webhomes.maths.ed.ac.uk/~v1ranick/papers/rota1.pdf}{Original paper}.
\bibitem{gem} D. Lopez-Paz and M. A. Ranzato. \emph{Gradient Episodic Memory
for Continual Learning}. 2017. \href{https://arxiv.org/abs/1706.08840}{Primary paper}.
\bibitem{more} J. J. Mor\'e and D. C. Sorensen. \emph{Computing a Trust Region
Step}. 1983. \href{https://epubs.siam.org/doi/10.1137/0904038}{Publisher article}.
\bibitem{elghaoui} L. El Ghaoui and H. Lebret. \emph{Robust Solutions to
Least-Squares Problems with Uncertain Data}. 1997.
\href{https://people.eecs.berkeley.edu/~elghaoui/Pubs/rob-ls.pdf}{Author-hosted paper}.
\bibitem{unitexture} \emph{UniTexture: Cross-Task Universal Adversarial Textures
for Vision-Language-Action Models}. 2026.
\href{https://arxiv.org/html/2608.13453v1}{Primary paper}.
\end{thebibliography}
\end{document}
